\begin{statementbox}
	\textbf{\textcolor{CStatement}{条件}}
	\begin{description}[style=nextline, leftmargin=2.8em, labelsep=0.5em, itemsep=0.35em]
		\item[\textbf{\textcolor{CStatement}{条件1（双曲线方程）：}}]
			双曲线标准方程为 $\frac{x^2}{a^2} - \frac{y^2}{b^2} = 1$（焦点在 $x$ 轴）或 $\frac{y^2}{a^2} - \frac{x^2}{b^2} = 1$（焦点在 $y$ 轴）。
		\item[\textbf{\textcolor{CStatement}{条件2（点P外部）：}}]
			点 $P(x_0,y_0)$ 在双曲线外部，且可引两条实切线。
	\end{description}
	\textbf{\textcolor{CStatement}{结论}}
	\begin{description}[style=nextline, leftmargin=2.8em, labelsep=0.5em, itemsep=0.45em]
		\item[\textbf{\textcolor{CStatement}{结论一（焦点x轴）：}}]
			\textbf{\textcolor{CStatement}{直观描述：}}
			焦点在 $x$ 轴时，切点弦所在直线方程形式为
			\[
				\frac{x_0 x}{a^2} - \frac{y_0 y}{b^2} = 1
			\]
		\item[\textbf{\textcolor{CStatement}{推广结论（焦点y轴）：}}]
			\textbf{\textcolor{CStatement}{直观描述：}}
			焦点在 $y$ 轴时，切点弦所在直线方程形式为
			\[
				\frac{y_0 y}{a^2} - \frac{x_0 x}{b^2} = 1
			\]
	\end{description}

\textbf{两有限实切点的存在条件。}
记 $q=x_0^2/a^2-y_0^2/b^2$，则恰有两个不同的有限实切点当且仅当
$q<1$ 且 $q\ne0$。$q=1$ 是一个切点，$q>1$ 无实切点；
$q=0$ 的非中心点在渐近线上，只有一个有限切点，另一方向为渐近线；中心点没有实切线。
纵轴形式使用 $q=y_0^2/a^2-x_0^2/b^2$。
\end{statementbox}
